% =====================================================================
% 02-kd-basics.tex -- Distillation basics.  T2, 8 min
% Sources: Hinton et al. 2015 and Furlanello et al. 2018 (seed papers),
%          Kim & Rush 2016; notes/emre-notes-on-sources.txt (Mansourian
%          survey summary). Theory note: notes/theory/t2-soft-targets-temperature.md
% =====================================================================
\renewcommand{\sectioninfo}{Presenter: T2 \quad $\cdot$ \quad 8 minutes}
\section{Distillation basics}

% ------------------------------------------------------------------
\begin{frame}{Soft targets carry more than the right answer\hfill\est\ \who{T2}}
  \begin{columns}[c]
    \begin{column}{0.46\textwidth}
      \small
      \begin{itemize}
        \item<1-> \textbf{Hard label}: only the correct next token counts
              (one-hot).
        \item<1-> \textbf{Soft target}: the teacher's probability for
              \emph{every} token, including the wrong ones.
        \item<2-> Which wrong answers look almost right is information:
              Hinton calls it \emph{dark knowledge}.
        \item<3-> A temperature $T$ softens the distribution and makes
              that information visible:
              \[
                p_i(T) = \frac{\exp(z_i/T)}{\sum_j \exp(z_j/T)}
              \]
      \end{itemize}
      \vspace{-0.15cm}
      \src{Hinton, Vinyals and Dean, 2015}
    \end{column}
    \begin{column}{0.52\textwidth}
      \centering
      \begin{tikzpicture}
        \begin{axis}[kdaxis,
            width=7.0cm,height=5.0cm,
            ybar,bar width=8pt,
            ymin=0,ymax=0.85,
            symbolic x coords={a,b,c,d,e},
            xtick=data,
            xticklabels={18,17,16,81,x},
            ylabel={probability},
            xlabel={next token after ``... so the answer is''},
            legend pos=north east,
            legend cell align=left,
            enlarge x limits=0.15]
          \addplot[fill=kdblue,draw=kdblue] coordinates
            {(a,0.745) (b,0.166) (c,0.061) (d,0.023) (e,0.005)};
          \addplot[fill=kdorange,draw=kdorange,visible on=<3->] coordinates
            {(a,0.393) (b,0.239) (c,0.171) (d,0.123) (e,0.074)};
          \legend{$T=1$, $T=3$}
        \end{axis}
      \end{tikzpicture}

      {\tiny\color{kdgray} Illustration with made-up logits $z=(4.0,\,2.5,\,1.5,\,0.5,\,-1.0)$.}
    \end{column}
  \end{columns}
\end{frame}

% ------------------------------------------------------------------
\begin{frame}{The distillation loss\hfill\est\ \who{T2}}
  \begin{columns}[T]
    \begin{column}{0.55\textwidth}
      \small
      Match the student's distribution $p_S$ to the teacher's $p_T$ with
      the Kullback--Leibler divergence, token by token:
      \[
        \mathcal{L} \;=\; T^2 \cdot \mathrm{KL}\!\left(p_T^{(T)} \,\middle\|\, p_S^{(T)}\right)
        \;=\; T^2 \sum_{v} p_T^{(T)}(v)\,\log\frac{p_T^{(T)}(v)}{p_S^{(T)}(v)}
      \]
      \begin{itemize}
        \item<2-> The $T^2$ factor keeps the gradient size comparable
              across temperatures \src{Hinton et al., 2015}.
        \item<3-> Classic KD often mixes in a hard-label term; we use the
              soft target alone, so the arms differ only in their text.
        \item<4-> \textbf{Our setting}: $T = 1$ for every arm (not swept;
              see limitations).
      \end{itemize}
    \end{column}
    \begin{column}{0.42\textwidth}
      \begin{block}{In our experiment}
        \footnotesize
        \arm{human} and \arm{ref\_hardlabel} train on the \emph{same}
        cleaned GSM8K solutions:
        \begin{itemize}
          \item \arm{human}: forward KL to the teacher's distribution
          \item \arm{ref\_hardlabel}: cross-entropy on the correct token
        \end{itemize}
        The difference between them measures what the soft targets add
        (results, part 7).
      \end{block}
    \end{column}
  \end{columns}
\end{frame}

% ------------------------------------------------------------------
\begin{frame}{Two seed papers, two ideas we build on\hfill\est\ \who{T2}}
  \begin{columns}[T]
    \begin{column}{0.48\textwidth}
      \begin{block}{Born-Again Networks \src{Furlanello et al., 2018}}
        \footnotesize
        A student with the \emph{same} architecture as its teacher, trained
        on the teacher's soft outputs, can end up \emph{better} than the
        teacher. Distillation is more than compression: the gain comes from
        the soft targets, not from the teacher being bigger.
      \end{block}
    \end{column}
    \begin{column}{0.48\textwidth}
      \begin{block}{Sequence-level KD \src{Kim and Rush, 2016}}
        \footnotesize
        Instead of matching the teacher token by token on fixed text, train
        the student on whole sequences the \emph{teacher} wrote. In practice:
        generate teacher answers once, fine-tune on them. This is our
        \arm{teacher} arm, and the DeepSeek-R1 recipe.
      \end{block}
    \end{column}
  \end{columns}
  \vspace{0.2cm}
  \begin{todobox}{T2's own examples (theory note is still empty)}
    Fill \repofile{notes/theory/t2-soft-targets-temperature.md} and add here, if
    useful: a real teacher distribution for one GSM8K token from our
    teacher, and a one-sentence take-away for the audience.
  \end{todobox}
\end{frame}
